%* auth-exec *%

\crefname{figure}{Fig.}{Figs.}
\Crefname{figure}{Fig.}{Figs.}
\crefname{table}{Tab.}{Tabs.}
\Crefname{table}{Tab.}{Tabs.}
\crefname{section}{Sect.}{Sects.}
\Crefname{section}{Sect.}{Sects.}
\crefname{req}{Req.\@}{Reqs.\@}
\Crefname{req}{Req.\@}{Reqs.\@}

\pgfplotsset{
    width=5cm, 
    legend style={
        font=\footnotesize,
        %rounded corners=2pt
    },
    compat=newest
}

\DeclareUnicodeCharacter{2212}{−}
\usepgfplotslibrary{groupplots,dateplot}
\usetikzlibrary{calligraphy, calc, arrows.meta, decorations.pathreplacing, matrix, positioning, fit, backgrounds, patterns,shapes.arrows, shapes.symbols}

\newlist{paraenum}{enumerate*}{1}
\setlist[paraenum]{label=\emph{(\arabic*)}}

\makeatletter
\@ifdefinable\SC@listings@vpos{\def\SC@listings@vpos{b}}
\newenvironment{SClisting}{\SC@float[\SC@listings@vpos]{mylisting}}{\endSC@float}
\newenvironment{SClisting*}{\SC@dblfloat[\SC@listings@vpos]{mylisting}}{\endSC@dblfloat}
%% `sidecap` and `float` access the name of the new float differently.
%% You can't rely on `float` to pass `sidecap` the correct name.  So,
%% here we manually feed the correct name for the new float in a manner 
%% pleasing to `sidecap`.
\@namedef{mylistingname}{Listing}
\makeatother

\lstdefinelanguage{RL}{
  keywords={on, if, then, else, outputs},
  ndkeywords={},
  sensitive=false,
  comment=[l]{//},
  columns=fullflexible
}

\lstdefinelanguage{Sancus-C}[]{C}{
    morekeywords={SM_ENTRY,SM_FUNC,SM_DATA,SM_INPUT,SM_OUTPUT},
}

\lstdefinelanguage{Reactive}{
  morekeywords={module, on, if, else},
  morecomment=[l]{\#}
}

\lstset{
  inputpath=listings,
  frame=lines,
  basicstyle=\ttfamily\tiny,
  keywordstyle=\bfseries\ttfamily,
  ndkeywordstyle=\bfseries\ttfamily,
  commentstyle=\itshape\ttfamily,
  stringstyle=\ttfamily,
  numberbychapter=false,
  showlines=true
}

%* v2x *%

\newtheorem{req}{Requirement}
\newtheorem{defn}{Definition}
\newtheorem{prop}{Property}

\lstdefinelanguage{Tamarin}{
  keywords={lemma,All,Ex,not,rule, restriction},
  keywordstyle = {\bfseries},
  otherkeywords={&,==>,.,=,<,>},
  comment=[l]{//},
  morecomment = [s]{/*}{*/},
  morestring=*[d]{'},
  columns=fullflexible
}

\lstdefinestyle{mystyle}{
    inputpath=listings,
    commentstyle=\color{teal},
    keywordstyle=\color{magenta},
    numberstyle=\tiny\color{gray},
    stringstyle=\color{purple},
    basicstyle=\ttfamily\scriptsize,
    breakatwhitespace=false,         
    breaklines=true,                 
    captionpos=b,                    
    keepspaces=true,                 
    numbers=none,                    
    numbersep=5pt,                  
    showspaces=false,                
    showstringspaces=false,
    showtabs=false,                  
    tabsize=2
}

\lstset{style=mystyle}

\hyphenation{pseu-do-nym}
\hyphenation{pseu-do-nyms}